% =====================================================================
%  FICHE 11 — THEME 4 — Exercices DIFFICILE — Potentiels standard et prévision
% =====================================================================
\documentclass[11pt,a4paper]{article}
\usepackage[utf8]{inputenc}
\usepackage[T1]{fontenc}
\usepackage{lmodern}
\usepackage[french]{babel}
\usepackage{amsmath,amssymb}
\usepackage[version=4]{mhchem}
\usepackage{siunitx}
\sisetup{output-decimal-marker={,},per-mode=symbol}
\usepackage{xcolor}
\usepackage{array}
\usepackage{booktabs}
\usepackage{enumitem}
\usepackage[most]{tcolorbox}
\usepackage[a4paper,margin=2.1cm,headheight=26pt,headsep=10pt]{geometry}
\usepackage{fancyhdr}
\usepackage[hidelinks]{hyperref}

\definecolor{primary}{RGB}{146,95,160}
\definecolor{primarydark}{RGB}{92,52,104}
\newcommand{\NO}{\ensuremath{\mathrm{N.O.}}}
\newcommand{\Ered}{\ensuremath{E_{\mathrm{r}}^{\circ}}}

\newcounter{exo}
\newtcolorbox{exo}[1][]{enhanced,breakable,boxrule=0.9pt,arc=2pt,colback=primary!4,
  colframe=primary,fonttitle=\bfseries,coltitle=white,left=3mm,right=3mm,top=2.2mm,bottom=2.2mm,
  title={Exercice~\refstepcounter{exo}\theexo\ifx\relax#1\relax\relax\else\ \textmd{--- \itshape #1}\fi}}
\newtcolorbox{donnees}{enhanced,boxrule=0.7pt,arc=2pt,colback=black!3,colframe=black!45,
  left=3mm,right=3mm,top=2mm,bottom=2mm,fonttitle=\bfseries\small,coltitle=black!70,
  title={Données — potentiels standard de réduction \Ered\ (à \SI{25}{\celsius})}}

\pagestyle{fancy}\fancyhf{}
\fancyhead[L]{\small\textcolor{primarydark}{\textbf{Fiche 11 · Thème 4} — Potentiels standard et prévision}}
\fancyhead[R]{\small\textcolor{primarydark}{\textsc{Difficile}}}
\fancyfoot[C]{\small\thepage}
\renewcommand{\headrulewidth}{0.8pt}
\renewcommand{\headrule}{\hbox to\headwidth{\color{primary}\leaders\hrule height \headrulewidth\hfill}}

\begin{document}
\begin{center}
{\LARGE\bfseries\color{primarydark}Thème 4 — Niveau Difficile}\\[2pt]
{\small\itshape Raisonner sur toute une échelle, prévoir des réactions concurrentes et des médiamutations.}\\[-4pt]
{\color{primary}\rule{\textwidth}{1pt}}
\end{center}

\begin{donnees}
\centering\small
\begin{tabular}{lccccc}
Couple & \ce{Cl2}/\ce{Cl^-} & \ce{Cr2O7^{2-}}/\ce{Cr^{3+}} & \ce{Br2}/\ce{Br^-} & \ce{Ag+}/\ce{Ag} & \ce{Fe^{3+}}/\ce{Fe^{2+}}\\
\Ered\ (V) & $+1{,}39$ & $+1{,}23$ & $+1{,}06$ & $+0{,}80$ & $+0{,}77$\\[3pt]
Couple & \ce{I2}/\ce{I^-} & \ce{Cu^{2+}}/\ce{Cu} & \ce{Sn^{4+}}/\ce{Sn^{2+}} & \ce{S4O6^{2-}}/\ce{S2O3^{2-}} & \ce{Zn^{2+}}/\ce{Zn}\\
\Ered\ (V) & $+0{,}536$ & $+0{,}34$ & $+0{,}151$ & $+0{,}08$ & $-0{,}76$\\
\end{tabular}
\end{donnees}

\begin{exo}[une lame de zinc dans une solution mixte]
On dispose des couples \ce{Zn^{2+}}/\ce{Zn} ($-0{,}76$), \ce{Cu^{2+}}/\ce{Cu} ($+0{,}34$) et
\ce{Ag+}/\ce{Ag} ($+0{,}80$).
\begin{enumerate}[label=\textbf{\alph*)},leftmargin=2.2em,topsep=2pt,itemsep=2pt]
  \item Classe ces trois oxydants par pouvoir oxydant croissant, et ces trois réducteurs par pouvoir
        réducteur croissant.
  \item Parmi \ce{Zn}, \ce{Cu}, \ce{Ag}, lesquels peuvent être oxydés par \ce{Cu^{2+}} ? Écris la (les) réaction(s).
  \item On plonge une lame de zinc dans une solution contenant \emph{à la fois} des ions \ce{Cu^{2+}}
        et \ce{Ag+}. Le zinc peut-il réduire ces deux ions ? Écris les deux réactions.
  \item Lequel des deux ions est réduit \textbf{en premier} ? Justifie par le pouvoir oxydant.
  \item Explique pourquoi un bijou en argent ne se dissout pas dans une solution de sulfate de cuivre.
\end{enumerate}
\end{exo}

\begin{exo}[une médiamutation du manganèse — d'après annales 2020]
On considère la réaction non pondérée :
\[
\ce{MnSO4 + KMnO4 + H2O -> MnO2 + H2SO4 + K2SO4}.
\]
\begin{enumerate}[label=\textbf{\alph*)},leftmargin=2.2em,topsep=2pt,itemsep=2pt]
  \item Détermine $\NO(\ce{Mn})$ dans \ce{KMnO4}, dans \ce{MnSO4} (ion \ce{Mn^{2+}}) et dans \ce{MnO2}.
  \item Le manganèse part de \emph{deux} degrés différents pour arriver à un \emph{seul} : identifie
        lequel est oxydé et lequel est réduit (c'est une \emph{médiamutation}).
  \item \ce{KMnO4} est-il l'oxydant ou le réducteur ? Un étudiant écrit : \og{}\ce{KMnO4} est l'oxydant
        car il \emph{cède} des électrons à \ce{MnSO4}\fg{}. Qu'y a-t-il d'erroné dans cette justification ?
  \item Écris la condition sur les potentiels standard pour que cette réaction soit spontanée
        (compare $\Ered(\ce{MnO4^-}/\ce{MnO2})$ et $\Ered(\ce{MnO2}/\ce{Mn^{2+}})$).
\end{enumerate}
\end{exo}

\begin{exo}[trois réactions : lesquelles ont lieu ?]
À l'aide du tableau de données, dis pour chacune si elle est \textbf{spontanée} dans les conditions
standard (calcule $\Delta E^{\circ}$ et, si elle n'a pas lieu, indique le sens spontané) :
\begin{enumerate}[label=\textbf{\alph*)},leftmargin=2.2em,topsep=2pt,itemsep=2pt]
  \item $\ce{2 I^- + S4O6^{2-} -> I2 + 2 S2O3^{2-}}$
  \item $\ce{Cr2O7^{2-} + 3 Sn^{2+} + 14 H+ -> 2 Cr^{3+} + 3 Sn^{4+} + 7 H2O}$
  \item $\ce{Br2 + 2 Cl^- -> 2 Br^- + Cl2}$
\end{enumerate}
\end{exo}

\end{document}
